\chapter*{Abstract}

\section*{\tituloPortadaEngVal}

Nowadays, many television programs, especially news broadcasts, make use of augmented reality (AR) elements displayed on screen to complement the information delivered by presenters. Thanks to this technology it is possible to integrate three-dimensional graphics, interactive data, or complex visualizations directly into the studio, making them appear as part of the real space. This not only completes and facilitates the understanding of the content, but also gives the broadcast a more dynamic and visually appealing look.

This Final Degree Project studies how to apply augmented reality to live broadcasts on platforms such as YouTube or other streaming services. To this end, OBS Studio is used as the main production tool, analysing how it can be adapted to generate effects similar to those used in professional television.

The work focuses especially on the case of online programming contests, where AR can be very helpful for presenting information more clearly, facilitating interaction with the audience, and improving the visual quality of the content. In this way, the goal is to demonstrate how this technology can enrich both the viewer's experience and that of the presenter themselves.

\section*{Keywords}

\noindent Augmented Reality, Live Streaming, OBS plugin, 3D Models, ArUco Markers, OpenCV, DOMjudge


